\section{Related Work}
\label{sec:pho_related_work}

This section reviews prior work related to object-centric force policies for rigid body transport, focusing on learning-based pushing, model-based pushing, and sim-to-real transfer under contact uncertainty.

\paragraph{Learning-based pushing and mobile manipulation.}
Reinforcement learning has been applied to planar pushing under contact-rich dynamics. \citet{ferrandis2023nonprehensile} propose RL with multimodal categorical exploration to better cover hybrid contact behaviors. \citet{bergmann2025precision} study precision-focused RL for pushing to improve goal accuracy. Other work explores sensing combinations such as vision and proprioception to improve pushing under varying object and contact properties~\citep{cong2022vision,wang2023learning}. For mobile manipulators and legged systems, \citet{dadiotis2025dynamic} present model-free constrained RL for object goal pushing with a mobile manipulator, reporting robustness across object properties. For multi-robot settings, \citet{tang2024collaborative} address collaborative planar pushing of polytopic objects with coordinated contact forces. Most learning approaches operate in robot-centric action spaces, with generalization addressed through architecture design, randomization, or adaptation. This motivates the force-space abstraction developed in this chapter.

\paragraph{Model-based pushing and control-theoretic approaches.}
Planar pushing has been studied extensively through analytical models and control-theoretic formulations. Early work by \citet{mason1986mechanics} and \citet{lynch1996stable} characterized the quasi-static regime of frictional contact during pushing. Limit-surface models~\citep{goyal1991planar} compactly represent frictional force-moment pairs during planar sliding. More recently, optimization-based methods have cast pushing as constrained optimal control, handling hybrid contact modes explicitly. \citet{hogan2018reactive,hogan2020reactive} propose hybrid MPC formulations for reactive planar non-prehensile manipulation, and contact-implicit trajectory optimization enables planners to reason directly over contact interactions~\citep{bui2025push}. Impedance control~\citep{hogan1985impedance} provides a foundational framework for interaction control, with whole-body impedance formulations extending compliance to mobile manipulators~\citep{dietrich2016wholebody}. These approaches are physically grounded and interpretable but often rely on accurate contact models and nontrivial online optimization, limiting their applicability when object properties are unknown.

\paragraph{Sim-to-real transfer under contact uncertainty.}
The reality gap in pushing arises from uncertain contact parameters, unmodeled frictional effects, and differences between simulated and physical interaction. Domain randomization~\citep{tobin2017domain,peng2018sim} addresses this by training over diverse simulated conditions so that real-world variation appears in-distribution at deployment. For non-prehensile manipulation, \citet{lowrey2018simtoreal} demonstrate sim-to-real transfer of learned pushing policies using carefully identified simulators and model ensembles. Data-driven contact models have been shown to outperform analytical counterparts in planar contact settings~\citep{fazeli2017learning}, and broader surveys characterize how modeling assumptions and uncertainty-handling strategies shape pushing performance~\citep{stuber2020survey}.

The method in this chapter adopts a complementary strategy: the policy is learned in object-centric force space, and domain randomization is applied to the object's physical parameters rather than to a robot's dynamics. Sim-to-real transfer requires only estimating the contact-resistance descriptor and mapping learned forces to robot behavior through compliant impedance control.

\paragraph{Gap addressed by this chapter.}
Rigid body transport under unknown mass and friction sits between model-based pushing and robot-centric learned control. Model-based methods provide physical structure but depend on accurate contact models; robot-centric policies can be robust in trained regimes but often transfer poorly across embodiments and physical parameters. The framework in this chapter resolves this tension by representing the policy in object-centric force space and conditioning it on an interaction descriptor, enabling generalization across physical parameters and zero-shot deployment on a legged mobile manipulator.
